\chapter{Before you start}
\label{ch:i-before}

\section{What you need}

\begin{center}
\begin{tabularx}{\linewidth}{@{}>{\bfseries}p{0.27\linewidth}X@{}}
  \toprule
  Item & Notes \\
  \midrule
  Jueying X30 Pro & charged, with its hard emergency stop working; the HIKMICRO PTZ camera and the
    Hertzino HZ-FA-371 mounted and powered \\
  Skydroid H12 & charged; or any Android phone or tablet (touch sticks) \\
  Mac (or Linux laptop) & macOS on Apple Silicon recommended, 16 GB RAM, 30 GB free disk;
    Homebrew, Git with your GitHub SSH key, the \cmd{gh} CLI logged in \\
  Network cable & optional: the robot's Debug Ethernet port, to keep the Mac online over Wi-Fi
    while it talks to the robot (Section~\ref{sec:i-apt}) \\
  USB-C cable & to install the app on the H12 with \cmd{adb} \\
  Access & the GitHub repositories (x30\_ros, X30Controller; NETRA is private); a NETRA license
    with the substation module (Chapter~\ref{ch:i-netra}) \\
  \bottomrule
\end{tabularx}
\end{center}

\section{Addresses and passwords}
\label{sec:i-addresses}

\begin{center}
\begin{tabularx}{\linewidth}{@{}>{\bfseries}p{0.3\linewidth}p{0.25\linewidth}X@{}}
  \toprule
  What & Value & Notes \\
  \midrule
  Robot Wi-Fi & \code{X30*****} & the serial number of the robot; password \code{12345678}
    (Deep Robotics' default) \\
  SSH to every host & user \code{ysc} & password \code{'} (one single quote, Deep Robotics'
    default) \\
  Navigation host & \code{192.168.1.106} & the stack, its ROS master \\
  Perception host & \code{192.168.1.105} & the cameras, the recorder \\
  Motion host & \code{192.168.1.103} & commands to UDP 43893, reports from it to UDP 43897 \\
  HIKMICRO PTZ camera & \code{192.168.1.64} & its admin password: \code{HIKMICRO_PASSWORD} \\
  Hertzino acoustic camera & \code{192.168.2.204} & its password: \code{HERTZINO_PASSWORD} \\
  H12 gateway & UDP \code{9870} on .106 & the app's token: \code{X30_TELEOP_TOKEN} \\
  NETRA & port \code{80} on the Mac & the robot's token: \code{X30_NETRA_TOKEN} \\
  \bottomrule
\end{tabularx}
\end{center}

\begin{warnbox}[Secrets]
The camera passwords and the two tokens live only in the private settings files of the hosts
(\file{/etc/x30/*.env}, readable by root and \code{ysc}). They are never launch arguments or
ROS parameters: \cmd{roslaunch} prints every parameter, and anyone on the robot's network can
read them from the master.
\end{warnbox}

\section{Time needed}

About an hour on the Mac (most of it downloads: NETRA's images and its AI model), half an hour on
the robot for the build, half an hour for the services, the app and NETRA, and an hour for the
checks of a new robot.
